%% ====================================================================
%%  "When Does Adaptive Recovery Pay?"  --  v3, IEEE conference format
%%
%%  SINGLE FILE. Compiles on its own with no image files present:
%%  every figure renders as a labelled placeholder box of the correct
%%  size, so the page count is already final at 6 pages.
%%
%%  When the figure PDFs exist in ./figs/ , flip ONE switch:
%%      \usefigsfalse   ->   \usefigstrue
%%  The boxes are replaced by the real images and nothing reflows,
%%  provided each figure is generated at the exact figsize given in
%%  the accompanying generation prompt.
%% ====================================================================
\documentclass[conference]{IEEEtran}
\IEEEoverridecommandlockouts

\usepackage{cite}
\usepackage{amsmath,amssymb}
\usepackage{graphicx}
\usepackage{xcolor}
\usepackage[T1]{fontenc}

%% -------- figure placeholder machinery -------------------------------
\newif\ifusefigs
\usefigstrue            % <-- set to \usefigstrue once figs/*.pdf exist

\newcommand{\figph}[3]{%   #1 = width, #2 = height, #3 = file
  \ifusefigs
    \includegraphics[width=#1]{#3}%
  \else
    \fbox{\parbox[c][#2][c]{\dimexpr#1-2\fboxsep-2\fboxrule\relax}{%
      \centering\footnotesize\color{gray}\texttt{\detokenize{#3}}\par
      \vspace{2pt}[figure placeholder]}}%
  \fi}
%% ---------------------------------------------------------------------

\begin{document}

\title{When Does Adaptive Recovery Pay?\\
Crossover-Horizon Prediction for Failure Recovery\\
under a Noisy Failure Classifier}

%% TODO: replace "Divina" with her full published name, add her e-mail,
%%       and confirm the department line.
\author{\IEEEauthorblockN{Pio Micheal}
\IEEEauthorblockA{\textit{Dept. of Artificial Intelligence and Data Science}\\
\textit{Karunya Institute of Technology and Sciences}\\
Coimbatore, India\\
piomicheal17@gmail.com}
\and
\IEEEauthorblockN{Divina}
\IEEEauthorblockA{\textit{Dept. of Data Science and Cyber Security}\\
\textit{Karunya Institute of Technology and Sciences}\\
Coimbatore, India\\
divina@karunya.edu}
}

\maketitle

\begin{abstract}
Tool-using agents fail in recurring and classifiable ways, and the choice of
recovery strategy is increasingly treated as something to be learned rather
than fixed in advance. A designer weighing a hand-written recovery table
against an online learner needs to know at what deployment horizon the learner
overtakes the table. Existing crossover analyses assume the policy observes the
true failure type, which real systems do not provide; the policy sees a
classifier's prediction. We reformulate the crossover problem under an
arbitrary confusion matrix by mapping the strategy-success matrix into
predicted-type space with an exact Bayes mixture, then separate the resulting
degradation into a reward-ceiling channel and a learnability channel. At a
classifier with 0.939 mean diagonal accuracy, ceiling loss exceeds learnability
loss by roughly six to one. The two channels also work against each other:
classifier noise raises the static table's regret rate by 58\% yet shifts the
predicted crossover only from 3053 to 3126 episodes, and we show by
construction that the horizon is not monotonic in classifier accuracy. The
reformulation's free constant was tested under a protocol registered before the
supporting data existed. The held-out ratio $R = 1.4729$ clears its
pre-committed threshold of 1.5, but the bootstrap interval $[1.00, 9.55]$
required by the same protocol does not confirm that the clearance is stable, so
we report the component as partially supported.
\end{abstract}

\begin{IEEEkeywords}
autonomous agents, failure recovery, multi-armed bandits, Thompson sampling,
noisy context, pre-registration, negative results
\end{IEEEkeywords}

\section{Introduction}

Language agents that call external tools fail often enough that failure
handling has become a design concern in its own right. Timeouts, tool
exceptions, malformed outputs and lost context dominate the failure inventories
of both deployed and benchmarked systems \cite{cemri,zhu,song}. The usual
answer is a fixed mapping $\pi$ from failure type to recovery strategy. Such
tables are cheap and easy to inspect, and can be quite good, but they never
improve; a table is exactly as good as its author's beliefs on the day it was
written.

The alternative is to treat each pair of failure type and recovery strategy as
a bandit arm and learn its value from outcomes. Adaptation then costs
exploration: early in a deployment the learner does worse than a good table
because it is still paying to find out what works, so whether adaptation is
worthwhile depends on how long the system will run. The designer wants a
number, the horizon at which the two curves cross.

A crossover predictor supplies that number. It computes the table's constant
per-episode regret rate exactly, models the learner's cumulative regret with an
asymptotic bound in the tradition of Lai and Robbins \cite{lai}, and solves for
the point where the two meet. The construction assumes the policy observes the
true failure type, and deployed systems do not supply one. Tracing the call
chain of the framework studied here confirmed that no policy can read a true
type; the context object has no such field, so every policy acts on a
classifier's prediction (Fig.~\ref{fig:pipeline}). This is a property of the
system rather than a modelling choice, and it means the unmodified predictor
answers a question about a system nobody has.

\begin{figure}[!t]
\centering
\figph{\columnwidth}{0.338\columnwidth}{figs/fig1_pipeline.pdf}
\caption{Where the classifier sits. Policies observe only the predicted type
$t'$ and never the true type $t$, which is the structural fact motivating
Section~\ref{sec:noise}.}
\label{fig:pipeline}
\end{figure}

This paper works through the resulting extension. We give an exact
reformulation of the crossover problem under an arbitrary confusion matrix,
using a Bayes mixture that carries the success matrix into predicted-type space
and reduces to the original construction when the classifier is perfect. We
decompose classifier-induced degradation into ceiling loss and gap shrinkage,
measure both at a trained classifier's empirical matrix, and show that the
predicted horizon is not monotonic in accuracy. We then evaluate the
extension's one free constant on a 26-environment dataset generated after the
evaluation protocol was committed. A companion study, pre-registered in the
same way, asks whether the sign of adaptive advantage under drift reduces to a
low-dimensional statistic; it does not, and we report that as a negative
result. What we do not claim is that the extension makes the underlying
asymptotic bound accurate: held-out absolute errors stay near 40--60\%, and the
evaluation asks only whether the reformulation adds error on top of what the
base bound already carries.

\section{Related Work}

Work on agent failure has moved from anecdote toward structured taxonomy,
whether organised by multi-agent failure mode \cite{cemri}, by single-agent
module with a debugger that isolates the earliest outcome-changing error
\cite{zhu}, or by the complementary route of improving the environment rather
than the agent \cite{song}. On the response side, Reflexion writes reflections
on failed attempts into an episodic buffer \cite{reflexion}, CRITIC calls
external tools to verify and revise the model's own output \cite{critic}, and
PALADIN trains recovery behaviour from systematically injected tool failures
\cite{paladin}.

Our concern is different. We propose no new self-correction mechanism, and ask
instead how a recovery-strategy selector should be evaluated against a static
baseline once selection is treated as a bandit problem with persistent outcome
memory. The underlying trade-off is classical: Lai and Robbins established the
logarithmic lower bound \cite{lai}, Auer et al. showed logarithmic regret is
achievable uniformly in time \cite{auer}, and Thompson sampling
\cite{thompson}, whose crossover behaviour we study, has a modern treatment in
Russo et al. \cite{russo}.

The specific difficulty here, a learner conditioning on a corrupted view of the
context, has an established literature. Kirschner and Krause study bandits in
which the environment supplies a context distribution while the realisation
stays hidden \cite{kirschner}; Kim et al. find that under noisy features the
Bayes-optimal hypothesis can sit far from the underlying realizability function
\cite{kim}; and latent bandits, where arm rewards are known conditional on an
unobserved discrete state, are the closest structural analogue to a failure
type the policy never sees \cite{maillard,hong}. Our mixture is a finite,
discrete specialisation of the same marginalisation idea and we claim no
novelty for it. Three things are ours: the corruption channel is a measured
confusion matrix rather than an assumed noise model; the quantity of interest
is a crossover horizon against a hand-written baseline, which is what a
practitioner acts on, not a regret rate; and we treat the free multiplicative
constant as an object to be calibrated and held-out tested rather than folded
into $O(\cdot)$.

\section{Setup}

\subsection{Environments and policies}

Failures come from $N = 7$ types and are recovered with one of $K = 6$
strategies. The types are \texttt{tool\_timeout}, \texttt{tool\_error},
\texttt{wrong\_tool\_selected}, \texttt{hallucinated\_output},
\texttt{malformed\_output}, \texttt{context\_loss} and
\texttt{ambiguous\_instruction}; the strategies are \texttt{retry},
\texttt{replan}, \texttt{switch\_tool}, \texttt{re\_ground}, \texttt{verify}
and \texttt{clarify}.

An environment is a pair $(P, \pi)$. The matrix $P \in [0,1]^{N \times K}$
gives the success probability $P[t,k]$ of strategy $k$ against true type $t$,
with conditionally Bernoulli rewards. The mapping $\pi$ is a hand-authored
table, and $p$ is the distribution over true types. Failures are injected into
simulated episodes, classified, and recovered by a bandit selector whose
outcomes persist across episodes.

One environment, which we call canonical, has a hand-authored $P$ and a
hand-authored $\pi$ agreeing with the row argmax of $P$ on six of seven types;
we call this quality 6/7, a competent baseline rather than a straw man. The
other 25 environments are sampled. Each row draws a uniform baseline in
$[0.15, 0.55]$, then one randomly chosen strategy in that row is raised to the
row maximum plus a fixed gap, capped at $0.95$; all results here use
gap $= 0.20$. Those 25 reuse the canonical table instead of receiving one
authored against their own $P$, so baseline quality across the sweep is
uncontrolled, and in two environments the reused table matches none of the
seven row argmaxes; the 6/7 figure describes the canonical environment only.
Section~\ref{sec:heldout} returns to this.

\subsection{The predictor being extended}

We restate the base construction so the paper stands on its own. Write
\begin{equation}
\mathrm{oracle}(P, p) = \sum_{t} p[t] \max_{k} P[t,k],
\label{eq:oracle}
\end{equation}
where $t \in \{1,\dots,N\}$ indexes the true failure type, $k \in
\{1,\dots,K\}$ indexes the recovery strategy, $p[t]$ is the probability that an
episode's failure is of type $t$, and $P[t,k]$ is the probability that strategy
$k$ recovers a type-$t$ failure; $\mathrm{oracle}(P,p)$ is thus the per-episode
success rate of a policy that always plays the best strategy for the true type.
Correspondingly,
\begin{equation}
\mathrm{policy}(P, \pi, p) = \sum_{t} p[t]\, P[t, \pi(t)],
\label{eq:policy}
\end{equation}
where $\pi(t) \in \{1,\dots,K\}$ is the strategy the hand-written table assigns
to type $t$, so $\mathrm{policy}(P,\pi,p)$ is the table's per-episode success
rate.

The table's per-episode regret rate $\rho = \mathrm{oracle}(P,p) -
\mathrm{policy}(P,\pi,p)$ is constant, so its cumulative regret is exactly
$\rho T$. The learner's cumulative regret is modelled by a Lai--Robbins-style
bound in KL form,
\begin{equation}
R_{\mathrm{TS}}(T) = c\,N \sum_{t} p[t]
\!\!\sum_{k\,:\,P[t,k] < P[t,k_t^{*}]}\!\!
\frac{\log\!\big(\max(T/N, 1)\big)}{\mathrm{KL}\big(P[t,k],\, P[t,k_t^{*}]\big)},
\label{eq:rts}
\end{equation}
where $T$ is the deployment horizon in episodes, $k_t^{*} = \arg\max_k P[t,k]$
is the best strategy for type $t$, the inner sum runs over the strictly
suboptimal strategies in row $t$, $\mathrm{KL}(a,b)$ is the Bernoulli
divergence clipped away from 0 and 1 for numerical stability, $N$ is the number
of types, and $c > 0$ is a free multiplicative constant. The crossover horizon
$T^{*}$ solves $R_{\mathrm{TS}}(T) = \rho T$. The constant $c$ has no
closed-form value for a given problem instance and must be fitted; under
confusion we write it $\kappa$. It is the single free parameter whose
transferability we test in Section~\ref{sec:heldout}, and everything downstream
of it is calibration rather than proof.

Let $\mathrm{confusion}[t,t'] = \Pr(\text{predicted } t' \mid \text{true } t)$.
The base predictor implicitly sets this to the identity matrix $I$. Two
questions follow. First, what is the right crossover prediction once the policy
conditions on $t'$? This turns out to have an exact answer on the static side
and only an approximate one on the adaptive side. Second, does a constant
fitted under one classifier condition transfer to another? If a $\kappa$ fitted
on perfect-classifier data predicts noisy held-out cells about as well as one
fitted directly on noisy data, the reformulation has absorbed the classifier
effect and the remaining error is inherited. If not, the reformulation is
adding error of its own. Separating these requires held-out data and a decision
rule fixed before looking.

\section{Crossover Under Classifier Noise}
\label{sec:noise}

\subsection{The effective mixture}

Marginalising over true types gives the predicted-type distribution
\begin{equation}
p'[t'] = \sum_{t} p[t]\, \mathrm{confusion}[t,t'],
\label{eq:pprime}
\end{equation}
where $t'$ indexes the type the classifier reports,
$\mathrm{confusion}[t,t']$ is the probability that a true type $t$ is reported
as $t'$, and $p'[t']$ is therefore the probability that an episode lands in
predicted bucket $t'$. Bayes' rule inside each bucket gives the effective
success probability of strategy $k$ given a prediction of $t'$:
\begin{equation}
P_{\mathrm{eff}}[t',k] = \frac{\sum_{t} p[t]\,\mathrm{confusion}[t,t']\,P[t,k]}
{p'[t']},
\label{eq:peff}
\end{equation}
where the numerator sums the contribution of every true type that can be
reported as $t'$, weighted by how often that happens, and the denominator
normalises by the bucket's total mass, so $P_{\mathrm{eff}}[t',k]$ is the
probability that strategy $k$ succeeds when the classifier says $t'$.

Both are exact and both collapse to $p' = p$, $P_{\mathrm{eff}} = P$ at
$\mathrm{confusion} = I$. Equation~\eqref{eq:peff} is the object the rest of
the extension is built on: given that the classifier said $t'$, how well does
strategy $k$ really work, averaged over what the type actually was.

\subsection{Static and adaptive sides}

The static side composes two exact functions,
\begin{equation}
\rho_{\mathrm{conf}} = \mathrm{oracle}(P,p) - \sum_{t}\sum_{t'}
p[t]\,\mathrm{confusion}[t,t']\,P[t,\pi(t')],
\label{eq:rhoconf}
\end{equation}
where $\pi(t')$ is the strategy the table fires on seeing the \emph{predicted}
type $t'$, while $P[t,\pi(t')]$ scores that strategy against the \emph{true}
type $t$; $\rho_{\mathrm{conf}}$ is thus the table's per-episode regret rate
once its input is corrupted, and it is exact rather than asymptotic. The
adaptive side reuses \eqref{eq:rts} unmodified, evaluated on the mixture,
\begin{equation}
R^{\mathrm{conf}}_{\mathrm{TS}}(T) =
R_{\mathrm{TS}}(P_{\mathrm{eff}}, T, N, p', \kappa),
\label{eq:rtsconf}
\end{equation}
where $P_{\mathrm{eff}}$ and $p'$ replace $P$ and $p$ from \eqref{eq:rts}, $N$
is unchanged, and $\kappa$ is the fitted constant playing the role of $c$. The
horizon under confusion solves $R^{\mathrm{conf}}_{\mathrm{TS}}(T) =
\rho_{\mathrm{conf}} T$.

This is its own routine rather than a wrapper on the base solver, because it
mixes index spaces: $\rho_{\mathrm{conf}}$ lives in true-type space while the
regret term lives in predicted-type space. Unit tests confirm that at
$\mathrm{confusion} = I$ the extension reproduces the base solver bit for bit,
so it is a strict generalisation rather than a parallel implementation that
happens to agree.

\subsection{Two channels, two checks}

Classifier noise does damage in two ways that need not move together. Merging
true types into indistinguishable buckets destroys information and lowers the
achievable ceiling whatever the policy does,
\begin{equation}
\Delta_{\mathrm{ceil}} = \mathrm{oracle}(P,p) -
\mathrm{oracle}(P_{\mathrm{eff}}, p'),
\label{eq:dceil}
\end{equation}
where the first term is the best achievable per-episode success rate with
perfect type information and the second is the best achievable rate once types
have been merged into predicted buckets, so $\Delta_{\mathrm{ceil}} \geq 0$
measures reward permanently lost to merging, with equality exactly at
$\mathrm{confusion} = I$. Merging also changes the gap between the best and
second-best action inside each bucket, which changes how hard the arm is to
identify. We report the mean shrinkage of that top-two gap across buckets,
computed as an unweighted mean over contributing true rows. The first is a
ceiling effect, the second a learnability effect.

Two checks guard the construction. The first asks whether a bucket's best
action under $P_{\mathrm{eff}}$ could be optimal for none of the true types
feeding it; it does not fire at our measured confusion, and a constructed
three-type example confirms it does fire when its structural condition holds.
The second concerns monotonicity: while $\rho_{\mathrm{conf}}$ and
$\Delta_{\mathrm{ceil}}$ are provably monotonic in classifier accuracy,
$T^{*}$ is not. A two-type, two-strategy unit test has a bucket whose best arm
changes identity as accuracy crosses a threshold, making its top-two gap
V-shaped and zero at the flip, whatever $\kappa$ is.

\subsection{Assumptions}

Five assumptions carry the extension and they are not equally safe. A1, the
predicted-type bucket structure, was verified by tracing the implementation;
A2, conditionally Bernoulli rewards, is inherited unchanged. A3, a stationary
confusion matrix, is optimistic wherever the classifier itself degrades. A4,
the asymptotic bound, holds as $T \to \infty$ rather than at finite time and
still carries an uncalibrated constant. A5, scalar accuracy standing in for a
full matrix, is a simplification we sidestep by validating against the
classifier's actual, very non-uniform matrix.

\section{Protocol}

The classifier combines a TF-IDF representation (word 1--3 grams, 5000
features) with a one-hot tool identifier, feeding a linear SVM with balanced
class weights; calibrated probabilities come from a separate logistic
regression. It is trained on 300 template-generated failure events from the
same synthetic injector the simulator uses, and the matrix in
Fig.~\ref{fig:confusion} was measured on a disjoint batch of 150 events per
type under a different seed. The 0.939 accuracy is held out with respect to
instances but not with respect to distribution, so the matrix should be read as
a plausible \emph{shape} of classifier error, concentrated and asymmetric with
one dominant confusable type, and not as an accuracy that would survive
production failure text.

\begin{figure}[!t]
\centering
\figph{\columnwidth}{0.62\columnwidth}{figs/fig2_confusion.pdf}
\caption{Empirical confusion matrix of the trained classifier,
$\Pr(\text{predicted} \mid \text{true})$, mean diagonal accuracy 0.939. Five of
seven rows are effectively identity rows.}
\label{fig:confusion}
\end{figure}

The evaluation protocol was written and committed before any of the supporting
data existed. Of the 25 gap-family environments, 8 went to TEST by a fixed-seed
draw and the one non-gap-family environment stayed in TRAIN unconditionally,
giving 18 TRAIN and 8 TEST; two environments already known to have poor errors
landed in TEST and we did not redraw. Fitting used bounded 1-D least squares on
mean squared log-space error with the KL form fixed. Branch A (transfer) fits
$\kappa_A$ on TRAIN perfect-classifier cells and applies it unchanged to TEST
noisy cells; Branch B (in-condition) fits $\kappa_B$ on TRAIN noisy cells and
applies it to the same TEST cells, giving the best achievable same-condition
benchmark. The decision rule was $R = E_A/E_B$ on median absolute relative
error, with $R \leq 1.5$ committed to Case 1, the effect is absorbed and
residual error inherited, and $R > 1.5$ to Case 2, the extension adds error of
its own, reported with no rescue attempt. Uncertainty came from an
environment-level bootstrap with 2000 resamples, and a constant landing within
one order of magnitude of the base predictor's own best fit, $[0.00045,
0.045]$, was to be called a calibrated predictor rather than a validated bound.
One specification, no second attempt.

\section{Results}

\subsection{Decomposition and static-side validation}

At the classifier's empirical matrix (Fig.~\ref{fig:confusion}) the ceiling
reduction is $\Delta_{\mathrm{ceil}} = 0.0120$ and the mean gap shrinkage
across the seven buckets is $0.0020$, so ceiling loss runs about six times
larger than learnability loss (Fig.~\ref{fig:decomp}a). Both are real and
neither is zero, but the bulk of the measured degradation is information lost
to merging rather than a harder arm-identification problem. The split is
specific to this confusion structure and should not be read as a general ratio.
The concentration is severe (Fig.~\ref{fig:decomp}b):
\texttt{malformed\_output} alone carries roughly $0.36$ of the off-diagonal
mass and \texttt{tool\_error} a further $0.067$, while the remaining five rows
are effectively identity rows, which is why the two losses are not spread
evenly across buckets. The argmax-inversion check fires in none of the seven.

We then checked the exact static component of \eqref{eq:rhoconf} against an
existing table of six noise levels with 50 simulated seeds each, where $\beta$
interpolates linearly between uniform-random type prediction and the measured
matrix, so $\beta = 0$ is a classifier that guesses and $\beta = 1$ is the real
one. That table carries no adaptive trajectory, so it validates the static side
alone and is never used as evidence for the adaptive theory. As
Fig.~\ref{fig:decomp}c shows, the median absolute error is 0.93 percentage
points and no level exceeds 1.5. This is the strongest validation result
anywhere in this work, which is about what one expects from closed-form
arithmetic rather than an asymptotic approximation.

\begin{figure}[!t]
\centering
\figph{\columnwidth}{0.62\columnwidth}{figs/fig3_decomposition.pdf}
\caption{Decomposition at the measured classifier and static-side validation.
(a) Ceiling loss is about six times larger than learnability loss. (b) The
error is concentrated: one row carries roughly $0.36$ of the off-diagonal mass
while five rows are identity rows. (c) The exact static form of
\eqref{eq:rhoconf} tracks simulation across six noise levels $\beta$, 50 seeds
each, median absolute error 0.93 percentage points.}
\label{fig:decomp}
\end{figure}

\subsection{The horizon itself}

Applying the fitted transfer constant to the canonical environment produces the
number the construction exists to give (Fig.~\ref{fig:crossover}). Under a
perfect classifier the predicted crossover falls at 3053 episodes; under the
measured matrix it falls at 3126.

The direction repays a moment's attention, because the two channels pull
against each other and the smaller one wins. Classifier noise raises the
table's per-episode regret rate from $\rho = 0.0214$ to $\rho_{\mathrm{conf}} =
0.0338$, a 58\% increase from a classifier that is right 94\% of the time, and
that alone would pull the crossover earlier: a baseline bleeding faster is
caught sooner. It does not, because the learner's problem also became harder.
Ceiling reduction and gap shrinkage enter \eqref{eq:rts} through inverse-KL
denominators sensitive to small changes in the top-two gap, and the net effect
moves the horizon about 2.4\% later. This near-cancellation belongs to this
environment and this matrix; Fig.~\ref{fig:nonmono} sweeps the same effect over
accuracy in a minimal two-type instance, where the horizon diverges where the
bucket's best arm changes identity. Given held-out errors of 40--60\%, the
honest resolution of these figures is ``roughly three thousand episodes'', not
3053 against 3126.

\begin{figure}[!t]
\centering
\figph{\columnwidth}{0.55\columnwidth}{figs/fig5_crossover.pdf}
\caption{The crossover on the canonical environment. Classifier noise steepens
the static baseline by 58\% but also flattens the learner's curve, so the
horizon moves 2.4\% later rather than earlier.}
\label{fig:crossover}
\end{figure}

\begin{figure}[!t]
\centering
\figph{\columnwidth}{0.56\columnwidth}{figs/fig6_nonmonotone.pdf}
\caption{A better classifier does not imply an earlier crossover. $T^{*}$
diverges at the accuracy where the bucket's best arm flips and its top-two gap
vanishes.}
\label{fig:nonmono}
\end{figure}

\subsection{Held-out evaluation of the constant}
\label{sec:heldout}

The power check cleared comfortably, with 18 of 18 TRAIN and 8 of 8 TEST cells
finite in both branches against a floor of 5. Fig.~\ref{fig:heldout} has the
headline numbers: Branch A reaches 59.7\% median absolute relative error
against Branch B's 40.5\%, for $R = 0.5970/0.4053 = 1.4729$.

The point estimate satisfies the committed rule, making this a Case 1 outcome:
the mixture captures the classifier-noise effect and what error remains is
inherited asymptotic-bound bias rather than something the extension introduced.
Two things temper that reading. Absolute errors of 59.7\% and 40.5\% are large,
so the test shows a pre-existing limitation transferring cleanly rather than
being resolved. And the margin is 1.8\%, with the pre-registered bootstrap
interval for $R$ running to 9.55, deep into Case 2, so a different draw of
held-out environments could plausibly have landed there; the fit-uncertainty
bootstrap agrees that the noisy-condition fit is markedly less stable than the
transferred one (Fig.~\ref{fig:heldout}a).

\begin{figure}[!t]
\centering
\figph{\columnwidth}{0.72\columnwidth}{figs/fig7_heldout.pdf}
\caption{Pre-registered held-out evaluation on the eight TEST cells.
(a) Fitted constants with 95\% bootstrap intervals; both sit inside the
pre-committed range and the in-condition fit is the less stable of the two.
(b) Branch A per-cell signed relative error, a spread no scalar constant can
absorb.}
\label{fig:heldout}
\end{figure}

Fig.~\ref{fig:heldout}b shows why the medians should not be read as a tight
summary: two environments over-predict by roughly $5.75\times$ under both
constants and one under-predicts by roughly $3\times$. No scalar constant fixes
this spread, and the protocol's no-rescue commitment meant we did not try to
introduce one. One candidate explanation deserves naming. The 25 sampled
environments reuse the canonical table, so baseline quality varies uncontrolled
and is very low in some cells; baseline quality sets $\rho$, and $\rho$ sets
the slope the regret curve must cross, so the heterogeneity plausibly tracks
baseline quality rather than anything about the reformulation. Controlling it,
or entering it as a pre-registered covariate, is the highest-priority fix to
this design. We flag it rather than analyse it now, because running a quality
covariate after seeing the per-cell errors is exactly the exploratory move the
protocol exists to prevent.

The fitted transfer constant, 0.0444, sits inside the pre-committed calibrated
range near its upper edge. We therefore call this component a calibrated
predictor with theoretical motivation rather than a validated bound, and
classify the crossover extension overall as partially supported. Case 1 was the
pre-committed outcome and is real held-out evidence for the first time in this
line of work, but the fragility is equally real. Equations \eqref{eq:pprime}--\eqref{eq:peff}, \eqref{eq:rhoconf}
and \eqref{eq:dceil}, the gap-shrinkage report, the argmax-inversion check and
the non-monotonicity finding are all exact and validated; only
\eqref{eq:rtsconf} and the horizon inherit the asymptotic approximation. All of
it is simulated, including the classifier, so the numbers describe the
construction rather than any live agent.

\subsection{Drift: a negative result}

A separate extension, independent of the first, asks whether a low-dimensional
statistic can predict the sign of adaptive advantage under environment drift,
out of sample. We fixed the falsifiable structure first: a candidate must beat
a majority-class guess under leave-one-cell-out evaluation. Three candidates
were pre-registered and tried in order over a pre-simulated grid of 12,800
cells across four drift mechanisms, with a fourth ruled out in advance if the
third failed. Fig.~\ref{fig:drift}a summarises the outcome. Survival and
gap-change was insufficient. Magnitude and recovery-ratio fitted well but was
retracted before it could count as a finding, because its recovery-ratio term
has the base predictor in its own denominator, so testing whether it predicts
adaptive advantage was largely testing whether the base bound has some validity
applied row by row. Environment regret and reversion reached $R^2 = 0.676$ but
its sign accuracy exactly equalled the majority baseline it had to beat. One
methodological wrinkle is worth recording: first-attempt statistics were means
over touched rows only, and most of magnitude's apparent contribution in the
second attempt ($+0.156$ $R^2$) was an artifact of that exclusion, shrinking to
$+0.006$ when recomputed at the environment level.

\begin{figure}[!t]
\centering
\figph{\columnwidth}{0.42\columnwidth}{figs/fig8_drift.pdf}
\caption{The drift study. (a) Three pre-registered candidate statistics: C1 is
insufficient, C2 was retracted for circularity, C3 fits well but its sign
accuracy only equals the majority-class baseline. (b) The one qualitative
finding that survives: adaptive advantage is negative exactly where staleness
damage reverts and positive where it does not.}
\label{fig:drift}
\end{figure}

One qualitative finding survives the failed quantitative rules: the sign of
adaptive advantage under drift is governed by whether staleness damage reverts,
not by how large it is (Fig.~\ref{fig:drift}b). All 13 negative-advantage cells
sit at a reversion fraction of at least 0.68 across 260 seed draws, and all 32
cells with reversion near zero come out positive across 640 draws. This also
explains why the static table beats the learner in 10 of 16 partial-reorder
cells: best and second-best swap on a period-2 cycle there, so the table's
stale action stays near-optimal while the learner pays exploration cost chasing
a target that returns to where it began. We treat this as correlational with a
mechanism behind it, and drift-type fixed effects retain $+0.092$ $R^2$ beyond
the two-statistic summary, structure we have not accounted for.

\section{Conclusion}

We extended a crossover-horizon predictor to the setting real systems occupy,
where recovery policies see a classifier's prediction rather than the true
failure type. The extension is exact where it can be, collapses to its
predecessor at a perfect classifier, and separates classifier-induced
degradation into a ceiling channel and a learnability channel, the former about
six times the latter at our measured accuracy. On the canonical environment the
two channels nearly cancel, leaving the horizon almost unmoved despite a 58\%
rise in the baseline's regret rate, and the horizon is not monotonic in
accuracy, which rules out an intuitive design heuristic. The free constant was
evaluated under a protocol fixed before the data existed: the criterion is
satisfied by point estimate but its uncertainty quantification does not confirm
that the satisfaction is stable, so we call the component a calibrated
predictor rather than a validated bound. The companion drift study stands as a
negative result with one qualitative finding intact.

Three directions follow, in order of priority. A live-agent evaluation should
measure $P$ and the confusion matrix on real tool failures, since every number
here rests on a synthetic template generator and an in-distribution classifier
accuracy that real failure logs would not be expected to reproduce. Static
baseline quality should then be controlled across the environment family, or
entered as a pre-registered covariate, so that the per-cell spread in
Fig.~\ref{fig:heldout}b can be attributed rather than merely flagged. Finally,
relaxing the stationary-confusion assumption would let the two extensions
connect, and a causal drift test with damage magnitude matched and reversion
varied alone would turn the surviving qualitative finding into a tested
mechanism.

\begin{thebibliography}{99}
\bibitem{cemri}
M. Cemri \emph{et al.}, ``Why do multi-agent LLM systems fail?,'' arXiv preprint arXiv:2503.13657, 2025.

\bibitem{zhu}
K. Zhu \emph{et al.}, ``Where LLM agents fail and how they can learn from failures,'' arXiv preprint arXiv:2509.25370, 2025.

\bibitem{song}
K. Song, A. Jayarajan, Y. Ding, Q. Su, Z. Zhu, S. Liu, and G. Pekhimenko, ``Aegis: Taxonomy and optimizations for overcoming agent-environment failures in LLM agents,'' arXiv preprint arXiv:2508.19504, 2025.

\bibitem{lai}
T. L. Lai and H. Robbins, ``Asymptotically efficient adaptive allocation rules,'' \emph{Advances in Applied Mathematics}, vol. 6, no. 1, pp. 4--22, 1985.

\bibitem{reflexion}
N. Shinn, F. Cassano, E. Berman, A. Gopinath, K. Narasimhan, and S. Yao, ``Reflexion: Language agents with verbal reinforcement learning,'' in \emph{Advances in Neural Information Processing Systems}, vol. 36, 2023, pp. 8634--8652.

\bibitem{critic}
Z. Gou, Z. Shao, Y. Gong, Y. Shen, Y. Yang, N. Duan, and W. Chen, ``CRITIC: Large language models can self-correct with tool-interactive critiquing,'' in \emph{Proc. Int. Conf. Learning Representations (ICLR)}, 2024.

\bibitem{paladin}
S. V. Vuddanti, A. Shah, S. K. Chittiprolu, T. Song, S. Dev, K. Zhu, and M. Chaudhary, ``PALADIN: Self-correcting language model agents to cure tool-failure cases,'' arXiv preprint arXiv:2509.25238, 2025.

\bibitem{auer}
P. Auer, N. Cesa-Bianchi, and P. Fischer, ``Finite-time analysis of the multiarmed bandit problem,'' \emph{Machine Learning}, vol. 47, no. 2--3, pp. 235--256, 2002.

\bibitem{thompson}
W. R. Thompson, ``On the likelihood that one unknown probability exceeds another in view of the evidence of two samples,'' \emph{Biometrika}, vol. 25, no. 3--4, pp. 285--294, 1933.

\bibitem{russo}
D. J. Russo, B. Van Roy, A. Kazerouni, I. Osband, and Z. Wen, ``A tutorial on Thompson sampling,'' \emph{Foundations and Trends in Machine Learning}, vol. 11, no. 1, pp. 1--96, 2018.

\bibitem{kirschner}
J. Kirschner and A. Krause, ``Stochastic bandits with context distributions,'' in \emph{Advances in Neural Information Processing Systems}, vol. 32, 2019, pp. 14090--14099.

\bibitem{kim}
J.-H. Kim, S.-Y. Yun, M. Jeong, J. Nam, J. Shin, and R. Combes, ``Contextual linear bandits under noisy features: Towards Bayesian oracles,'' in \emph{Proc. 26th Int. Conf. Artificial Intelligence and Statistics (AISTATS)}, PMLR, vol. 206, 2023, pp. 1624--1645.

\bibitem{maillard}
O.-A. Maillard and S. Mannor, ``Latent bandits,'' in \emph{Proc. 31st Int. Conf. Machine Learning (ICML)}, 2014, pp. 136--144.

\bibitem{hong}
J. Hong, B. Kveton, M. Zaheer, Y. Chow, A. Ahmed, and C. Boutilier, ``Latent bandits revisited,'' in \emph{Advances in Neural Information Processing Systems}, vol. 33, 2020, pp. 13423--13433.
\end{thebibliography}

\end{document}
